\chapter{Configuration and persistence}
\label{ch:config}

\begin{aside}
Every non-trivial app reads configuration from somewhere ---
defaults, environment variables, command-line flags, a config
file, a database. Persisting state across runs follows the same
pattern in reverse. \maya{} does not impose a config system, but
the architecture invites a specific pattern: read at startup,
hold in the model, never re-read inside the loop.
\end{aside}

\section{The startup pipeline}

A typical \maya{} app does the following before \code{run}:

\begin{enumerate}[leftmargin=*]
  \item Parse command-line arguments.
  \item Load the config file.
  \item Read relevant environment variables.
  \item Apply precedence: CLI overrides env overrides file
    overrides defaults.
  \item Construct the initial \code{Model} including the merged
    config.
  \item Hand the model to \code{run}.
\end{enumerate}

The startup pipeline is the only place that touches the
filesystem for config. Inside the loop, the model carries
everything.

\subsection*{Why not reload on demand?}

A frequent temptation: ``read the config file every few seconds
so changes take effect''. The Elm-style approach: model the
file as state, watch it with a \code{Cmd::task}, and turn
changes into messages. The runtime is the same; the contract
stays clean.

\section{Defaults}

A small struct with sensible defaults:

\begin{lstlisting}
struct Config {
    bool   show_help      = true;
    int    history_size   = 100;
    Theme  theme          = themes::Dark();
    std::string log_path  = \"~/.cache/myapp/log\";
    std::chrono::milliseconds idle_timeout{5000};
};

constexpr Config default_config{};
\end{lstlisting}

The struct holds everything user-tunable. Designated init
allows partial override:

\begin{lstlisting}
Config cfg{
    .theme = themes::Light(),
    .history_size = 500,
};
\end{lstlisting}

Fields not mentioned retain their defaults.

\section{Config files}

\subsection*{Format}

TOML, YAML, JSON, INI --- pick one. \maya{} examples use TOML
because:

\begin{itemize}[leftmargin=*]
  \item Human-readable.
  \item Strict parser available (\code{toml++}).
  \item No indentation sensitivity (unlike YAML).
  \item Comments allowed.
\end{itemize}

\subsection*{Loading}

\begin{lstlisting}
Expected<Config, std::string> load_config(std::string_view path) {
    return file::read(path)
        .transform_error([](auto& e) { return e.context; })
        .and_then(parse_toml)
        .transform(toml_to_config);
}
\end{lstlisting}

A monadic chain (Chapter~\ref{ch:expected}): read, parse,
convert. Each step can fail; the chain handles propagation.

\subsection*{Default location}

XDG conventions on Linux:

\begin{lstlisting}
std::string config_path() {
    if (auto* x = std::getenv(\"XDG_CONFIG_HOME\")) {
        return std::string{x} + \"/myapp/config.toml\";
    }
    return std::string{std::getenv(\"HOME\")}
        + \"/.config/myapp/config.toml\";
}
\end{lstlisting}

On macOS, \code{\$HOME/Library/Application Support/myapp/}; on
Windows, \code{\%APPDATA\%\textbackslash myapp\textbackslash}.
The platform abstraction layer (Chapter~\ref{ch:platform}) can
provide the helper.

\section{Environment variables}

For machine-specific overrides, env vars are common:

\begin{lstlisting}
if (auto* t = std::getenv(\"MYAPP_THEME\")) {
    cfg.theme = themes::by_name(t).value_or(cfg.theme);
}
if (auto* l = std::getenv(\"MYAPP_LOG\")) {
    cfg.log_path = l;
}
\end{lstlisting}

Convention: prefix env vars with your app name to avoid
collisions.

\subsection*{Standard variables to honour}

A small list of variables most users expect:

\begin{itemize}[leftmargin=*]
  \item \code{NO\_COLOR}: disable colour output.
  \item \code{FORCE\_COLOR}: enable colour even on non-tty.
  \item \code{COLORTERM}: terminal colour capabilities.
  \item \code{TERM}: terminal type.
  \item \code{HOME}, \code{XDG\_CONFIG\_HOME}, etc.: paths.
  \item \code{LANG}, \code{LC\_*}: locale (affects width tables).
\end{itemize}

\section{Command-line flags}

The simplest parser is hand-written; more elaborate apps use
\code{argparse} or \code{cxxopts}:

\begin{lstlisting}
int main(int argc, char** argv) {
    Config cfg = default_config;
    for (int i = 1; i < argc; ++i) {
        std::string_view arg = argv[i];
        if (arg == \"--theme\" && i + 1 < argc) {
            cfg.theme = themes::by_name(argv[++i]).value_or(cfg.theme);
        } else if (arg == \"--help\") {
            print_help();
            return 0;
        }
    }
    return maya::run<App>(cfg);
}
\end{lstlisting}

For complex apps with subcommands, use a library. The patterns
are the same.

\section{Precedence}

CLI overrides env overrides file overrides defaults. In code:

\begin{lstlisting}
Config build_config(int argc, char** argv) {
    Config cfg = default_config;
    if (auto file_cfg = load_config(config_path())) {
        merge(cfg, file_cfg.value());
    }
    apply_env(cfg);
    apply_cli(cfg, argc, argv);
    return cfg;
}
\end{lstlisting}

Each step layered on top. \code{merge} is field-by-field
override; only fields explicitly set in the higher layer
replace the lower layer.

\section{Persisting state}

The flip side of config: state that survives across runs.
Examples:

\begin{itemize}[leftmargin=*]
  \item Window size (the app should reopen at the same size).
  \item Last opened file.
  \item Cursor position.
  \item Command history.
\end{itemize}

The pattern:

\begin{enumerate}[leftmargin=*]
  \item On startup, load a state file (often JSON or binary).
  \item Restore the state into the model.
  \item On meaningful changes, write the state back (debounced
    to avoid thrashing the filesystem).
  \item On quit, write a final state.
\end{enumerate}

\subsection*{Write strategy: debounced}

\begin{lstlisting}
[&](StateChanged) {
    m.dirty = true;
    return std::pair{m, Cmd<Msg>::after(2s, FlushState{})};
}

[&](FlushState) {
    if (!m.dirty) return std::pair{m, Cmd<Msg>{}};
    m.dirty = false;
    return std::pair{m, save_state_cmd(m.persistable_state())};
}
\end{lstlisting}

A change sets a flag and schedules a flush. The flush, when it
fires, checks the flag and writes. Multiple changes within the
window collapse to one write.

\subsection*{Write atomically}

To avoid corrupted state files (if the app crashes mid-write):

\begin{lstlisting}
Expected<void, IoError> atomic_write(std::string_view path, std::string_view data) {
    std::string tmp = std::string{path} + \".tmp\";
    auto w = file::write(tmp, data);
    if (!w) return std::unexpect, w.error();
    return file::rename(tmp, path);
}
\end{lstlisting}

Write to a temp file, then atomically rename. The rename is
atomic on POSIX; on Win32 you need \code{ReplaceFileW}.

\section{Schema migration}

State files persist across versions. If you add a field or
change a type, old state files must still load. Two strategies:

\subsection*{Versioned schemas}

\begin{lstlisting}
struct StateV1 { int cursor; };
struct StateV2 { int cursor; std::string filename; };

StateV2 migrate(StateV1 s) {
    return StateV2{.cursor = s.cursor, .filename = \"\"};
}
\end{lstlisting}

Read the version field; pick the right struct; migrate forward.

\subsection*{Tolerant readers}

Each field is optional; missing fields use defaults:

\begin{lstlisting}
struct State {
    int cursor = 0;
    std::string filename = \"\";
    // ... older versions had fewer fields
};
\end{lstlisting}

The deserialiser accepts any TOML object and populates known
fields. Unknown fields are ignored (or logged).

\section{Secrets}

Some config values are sensitive (API keys, tokens). The
patterns:

\begin{itemize}[leftmargin=*]
  \item Read from a file that only the user can read
    (mode 0600).
  \item Read from a system keychain (libsecret, Keychain
    Services).
  \item Read from an environment variable that the user sets
    in their shell config.
\end{itemize}

Do not write secrets to the state file. If the user wants to
persist them, that is a separate, deliberate choice.

\section{Pitfalls}

\begin{warning}
\textbf{Reading config inside view or update.} The config
should be loaded once at startup and stored in the model.
Re-reading mid-loop is slow and racy.
\end{warning}

\begin{warning}
\textbf{Tight coupling to file format.} If your model is a
TOML-shaped struct, swapping to JSON requires changing the
model. Keep the on-disk shape separate from the in-memory
model.
\end{warning}

\begin{warning}
\textbf{Saving on every keystroke.} The filesystem will hate
you. Debounce; consider write barriers.
\end{warning}

\begin{warning}
\textbf{Forgetting to validate.} A config file written by a
human can contain garbage. Validate on load; reject with a
clear error.
\end{warning}

\section*{Workshop}

\begin{enumerate}[leftmargin=*]
  \item Write a small config struct with five fields. Load
    defaults; override one from an env var; verify in the
    model.
  \item Add atomic writes for a state file. Force-kill the
    app mid-write; verify the file is not corrupted.
  \item Implement schema migration from a V1 state file to a
    V2. Write a unit test for the migration.
\end{enumerate}

\begin{recap}
Read config once at startup; carry it in the model; never re-
read inside the loop. Files in TOML or JSON, environment
variables for overrides, CLI flags for ad-hoc tweaks. Persist
state debounced and atomically. Migrate schemas across
versions. Secrets stay out of state files.
\end{recap}
